\section{Anexos}\label{cap:anexos}

Los anexos contienen los instrumentos de recolección de datos y los formatos de consentimiento. Todos se aplicarán después de su validación por juicio de expertos. Los anexos A a G corresponden a los instrumentos de recolección de datos y a los formatos del estudio; los anexos H a Q reúnen la documentación técnica y de evaluación de la ampliación Lush Beauty Smart DSS (requisitos, historias de usuario, casos de prueba, diseño de API y de flujos n8n, instrumentos adicionales, lista de preparación, fichas de decisión y evaluación de decisiones).

\anexotabla{A}

\anexo{Anexo A. Cuestionario de requerimientos}

Cuestionario de requerimientos del sistema de gestión del salón Lush Beauty. Está dirigido a la administradora y a las especialistas, con una duración estimada de 25 minutos. No existen respuestas correctas; se solicita describir la operación actual del salón con la mayor precisión posible.


\textbf{Instrucciones:} marque con una X las opciones indicadas o escriba en los espacios previstos. Sus respuestas son confidenciales y se identifican con un código de participante. En el Anexo G se presenta un ejemplo de llenado.


Código de participante: \blank{8}     Fecha: \blank{4}/\blank{4}/\blank{8}


\subanexo{I. Datos generales}

\pregunta{\textbf{1.} Su rol en el salón:}

\opciones{[~~] Administración\quad{}[~~] Especialista\quad{}[~~] Ambos}

\pregunta{\textbf{2.} ¿Cuántos años lleva operando el salón? \blank{12}}

\pregunta{\textbf{3.} ¿Cuántas personas prestan servicios actualmente? \blank{12}}

\pregunta{\textbf{4.} ¿Cuántos servicios se realizan en un día habitual? \blank{12}}

\pregunta{\textbf{5.} ¿Con qué dispositivos se registran hoy los servicios?}

\opciones{[~~] Celular\quad{}[~~] Computadora\quad{}[~~] Papel\quad{}[~~] Hoja de cálculo\quad{}[~~] Otro: \blank{8}}

\subanexo{II. Registro de servicios y caja}

\pregunta{\textbf{6.} ¿Cómo registra hoy los servicios realizados?}

\opciones{[~~] Libreta\quad{}[~~] Hoja de cálculo\quad{}[~~] Mensajes\quad{}[~~] Otro: \blank{8}}

\pregunta{\textbf{7.} ¿Cuándo registra un servicio?}

\opciones{[~~] En el momento\quad{}[~~] Al final del día\quad{}[~~] Al final de la semana}

\pregunta{\textbf{8.} Marque los datos que registra por cada servicio:}

\opciones{[~~] Fecha\quad{}[~~] Hora\quad{}[~~] Cliente\quad{}[~~] Especialista\quad{}[~~] Servicio\quad{}[~~] Precio\quad{}[~~] Descuento\quad{}[~~] Método de pago\quad{}[~~] Propina\quad{}[~~] Observaciones}

\pregunta{\textbf{9.} ¿Qué métodos de pago acepta?}

\opciones{[~~] Efectivo\quad{}[~~] QR\quad{}[~~] Transferencia\quad{}[~~] Tarjeta\quad{}[~~] Otro: \blank{8}}

\pregunta{\textbf{10.} ¿Un cliente puede pagar varios servicios en un mismo cobro (combos)?}

\opciones{[~~] Sí, frecuentemente\quad{}[~~] Sí, a veces\quad{}[~~] No}

\pregunta{\textbf{11.} ¿Cómo cuadra la caja al cierre del día? \blank{46}}

\pregunta{\textbf{12.} ¿Qué descuentos aplica y quién los autoriza?}

\opciones{[~~] Porcentaje fijo\quad{}[~~] Monto fijo\quad{}[~~] Libre}

\pregunta{\textbf{13.} Cuando un servicio se cancela, ¿qué sucede con el registro y con la comisión? \blank{30}}

\subanexo{III. Comisiones y liquidación}

\pregunta{\textbf{14.} ¿Qué porcentaje de comisión recibe cada especialista? \blank{12}}

\pregunta{\textbf{15.} ¿Sobre qué monto se calcula la comisión?}

\opciones{[~~] Precio de lista\quad{}[~~] Monto cobrado\quad{}[~~] Otro: \blank{8}}

\pregunta{\textbf{16.} ¿La comisión cambia según el servicio, la categoría o la persona?}

\opciones{[~~] Sí\quad{}[~~] No}

\pregunta{\textbf{17.} Explique la respuesta anterior, si corresponde: \blank{46}}

\pregunta{\textbf{18.} ¿Cada cuánto tiempo se paga al personal?}

\opciones{[~~] Semanal\quad{}[~~] Quincenal\quad{}[~~] Mensual\quad{}[~~] Otro: \blank{8}}

\pregunta{\textbf{19.} ¿Existe un sueldo fijo además de la comisión?}

\opciones{[~~] Sí, monto: \blank{8}\quad{}[~~] No}

\pregunta{\textbf{20.} ¿Qué deducciones se aplican al personal, y cómo se calcula cada una?}

\opciones{[~~] Adelantos\quad{}[~~] Faltas\quad{}[~~] Materiales\quad{}[~~] Otras: \blank{8}}

\pregunta{\textbf{21.} ¿Qué sucede cuando las deducciones superan lo ganado en el periodo? \blank{30}}

\pregunta{\textbf{22.} ¿Cómo se entrega el pago y su detalle al personal?}

\opciones{[~~] Efectivo sin detalle\quad{}[~~] Efectivo con detalle impreso\quad{}[~~] Mensaje\quad{}[~~] Otro: \blank{8}}

\subanexo{IV. Propinas}

\pregunta{\textbf{23.} ¿Cómo se registran hoy las propinas? \blank{30}}

\pregunta{\textbf{24.} ¿Cómo se reparten las propinas?}

\opciones{[~~] 100~\% a quien atendió\quad{}[~~] Entre el equipo\quad{}[~~] Con retención del salón\quad{}[~~] Otro: \blank{8}}

\pregunta{\textbf{25.} ¿Necesita modificar una propina después de registrada?}

\opciones{[~~] Sí\quad{}[~~] No}

\subanexo{V. Clientes}

\pregunta{\textbf{26.} ¿Cómo lleva el registro de clientes?}

\opciones{[~~] No lo lleva\quad{}[~~] Libreta\quad{}[~~] Hoja de cálculo\quad{}[~~] Otro: \blank{8}}

\pregunta{\textbf{27.} ¿Qué datos necesita de cada clienta?}

\opciones{[~~] Nombre\quad{}[~~] Teléfono\quad{}[~~] Fecha de nacimiento\quad{}[~~] Preferencias\quad{}[~~] Notas\quad{}[~~] Otros: \blank{8}}

\pregunta{\textbf{28.} ¿Necesita ver el historial de servicios de una clienta?}

\opciones{[~~] Sí\quad{}[~~] No}

\pregunta{\textbf{29.} ¿Ha tenido dificultades con nombres duplicados o confusos? Describa: \blank{30}}

\subanexo{VI. Catálogo y precios}

\pregunta{\textbf{30.} ¿Con qué frecuencia cambian los precios? \blank{12}}

\pregunta{\textbf{31.} ¿Hay servicios cuyo precio varía al momento del cobro (por ejemplo, diseños)? Cuáles: \blank{30}}

\opciones{[~~] Sí\quad{}[~~] No}

\pregunta{\textbf{32.} ¿Quién debe poder modificar el catálogo y los precios? \blank{12}}

\subanexo{VII. Reportes, importación y exportación}

\pregunta{\textbf{33.} ¿Qué reportes necesita?}

\opciones{[~~] Caja diaria\quad{}[~~] Detalle por servicio\quad{}[~~] Liquidación por persona\quad{}[~~] Otros: \blank{8}}

\pregunta{\textbf{34.} ¿En qué formato los necesita?}

\opciones{[~~] Excel\quad{}[~~] PDF\quad{}[~~] Pantalla\quad{}[~~] Otro: \blank{8}}

\pregunta{\textbf{35.} ¿Necesita trasladar a la nueva herramienta la información de su planilla actual?}

\opciones{[~~] Sí\quad{}[~~] No}

\pregunta{\textbf{36.} ¿Quién necesita recibir reportes fuera del salón?}

\opciones{[~~] Contador\quad{}[~~] Familia o socios\quad{}[~~] Ninguno\quad{}[~~] Otro: \blank{8}}

\subanexo{VIII. Uso, conectividad y seguridad}

\pregunta{\textbf{37.} ¿Con qué frecuencia se queda sin internet mientras atiende?}

\opciones{[~~] Nunca\quad{}[~~] Pocas veces al mes\quad{}[~~] Casi a diario}

\pregunta{\textbf{38.} ¿Necesita registrar servicios cuando no hay internet?}

\opciones{[~~] Sí\quad{}[~~] No}

\pregunta{\textbf{39.} ¿Quién debe poder ver los sueldos y las comisiones? \blank{12}}

\pregunta{\textbf{40.} ¿Quién debe poder modificar o eliminar un registro? \blank{12}}

\pregunta{\textbf{41.} ¿Necesita saber quién cambió un dato, cuándo y cuál era el valor anterior?}

\opciones{[~~] Sí\quad{}[~~] No}

\subanexo{IX. Priorización de necesidades}

\pregunta{\textbf{42.} Califique cada necesidad del 1 (no es importante) al 5 (imprescindible), marcando una casilla por fila.}

\vspace{8pt}{\footnotesize\setstretch{1.0}
\begin{longtable}{>{\raggedright\arraybackslash}p{0.6251\dimexpr\linewidth-12\tabcolsep\relax}>{\raggedright\arraybackslash}p{0.0750\dimexpr\linewidth-12\tabcolsep\relax}>{\raggedright\arraybackslash}p{0.0750\dimexpr\linewidth-12\tabcolsep\relax}>{\raggedright\arraybackslash}p{0.0750\dimexpr\linewidth-12\tabcolsep\relax}>{\raggedright\arraybackslash}p{0.0750\dimexpr\linewidth-12\tabcolsep\relax}>{\raggedright\arraybackslash}p{0.0750\dimexpr\linewidth-12\tabcolsep\relax}}
\caption{Priorización de necesidades del salón}\label{tab:v4_A1}\\
\toprule
\textbf{Necesidad} & \textbf{1} & \textbf{2} & \textbf{3} & \textbf{4} & \textbf{5} \\
\midrule\endfirsthead
\toprule
\textbf{Necesidad} & \textbf{1} & \textbf{2} & \textbf{3} & \textbf{4} & \textbf{5} \\
\midrule\endhead
\bottomrule\endfoot
\bottomrule\endlastfoot
Registrar un servicio en menos de 15 segundos desde el celular &  &  &  &  &  \\
Calcular automáticamente la comisión de cada línea &  &  &  &  &  \\
Liquidación semanal con deducciones y excedentes &  &  &  &  &  \\
Comprobante en PDF para compartir por WhatsApp &  &  &  &  &  \\
Registrar servicios sin conexión a internet &  &  &  &  &  \\
Importar la planilla actual de Excel &  &  &  &  &  \\
Bitácora de cambios consultable &  &  &  &  &  \\
Reportes en Excel para el contador &  &  &  &  &  \\
Ficha de clientes con historial &  &  &  &  &  \\
Agenda de citas en línea (posible etapa posterior) &  &  &  &  &  \\
\end{longtable}}

\subanexo{X. Observaciones}

\pregunta{\textbf{43.} ¿Cuál es el problema que más le preocupa hoy en la gestión del salón? \blank{30}}

\pregunta{\textbf{44.} ¿Qué parte del proceso actual no debería cambiar? \blank{30}}

Gracias por su participación.


\anexotabla{B}

\anexo{Anexo B. Cuestionario de usabilidad (SUS, versión en español)}

Versión en español del \textit{System Usability Scale} (Brooke, 1996). La traducción es propia y debe validarse por juicio de expertos antes de su aplicación. Se responde después de usar el sistema durante el piloto. Escala: 1 = totalmente en desacuerdo; 5 = totalmente de acuerdo.


\pregunta{\textbf{1.} Creo que me gustaría usar este sistema con frecuencia.}

\opciones{[~~] 1\quad{}[~~] 2\quad{}[~~] 3\quad{}[~~] 4\quad{}[~~] 5}

\pregunta{\textbf{2.} Encontré el sistema innecesariamente complejo.}

\opciones{[~~] 1\quad{}[~~] 2\quad{}[~~] 3\quad{}[~~] 4\quad{}[~~] 5}

\pregunta{\textbf{3.} Pensé que el sistema era fácil de usar.}

\opciones{[~~] 1\quad{}[~~] 2\quad{}[~~] 3\quad{}[~~] 4\quad{}[~~] 5}

\pregunta{\textbf{4.} Creo que necesitaría el apoyo de una persona técnica para poder usar este sistema.}

\opciones{[~~] 1\quad{}[~~] 2\quad{}[~~] 3\quad{}[~~] 4\quad{}[~~] 5}

\pregunta{\textbf{5.} Encontré que las distintas funciones del sistema estaban bien integradas.}

\opciones{[~~] 1\quad{}[~~] 2\quad{}[~~] 3\quad{}[~~] 4\quad{}[~~] 5}

\pregunta{\textbf{6.} Pensé que había demasiadas inconsistencias en este sistema.}

\opciones{[~~] 1\quad{}[~~] 2\quad{}[~~] 3\quad{}[~~] 4\quad{}[~~] 5}

\pregunta{\textbf{7.} Imagino que la mayoría de las personas aprendería a usar este sistema muy rápidamente.}

\opciones{[~~] 1\quad{}[~~] 2\quad{}[~~] 3\quad{}[~~] 4\quad{}[~~] 5}

\pregunta{\textbf{8.} Encontré el sistema muy incómodo de usar.}

\opciones{[~~] 1\quad{}[~~] 2\quad{}[~~] 3\quad{}[~~] 4\quad{}[~~] 5}

\pregunta{\textbf{9.} Me sentí muy seguro al usar el sistema.}

\opciones{[~~] 1\quad{}[~~] 2\quad{}[~~] 3\quad{}[~~] 4\quad{}[~~] 5}

\pregunta{\textbf{10.} Necesité aprender muchas cosas antes de poder empezar a usar este sistema.}

\opciones{[~~] 1\quad{}[~~] 2\quad{}[~~] 3\quad{}[~~] 4\quad{}[~~] 5}

\textbf{Cálculo de la puntuación.} En los ítems impares se resta 1 a la respuesta; en los pares, la respuesta se resta de 5. La suma de los valores obtenidos se multiplica por 2,5, lo que da un resultado entre 0 y 100.


\anexotabla{C}

\anexo{Anexo C. Fichas de observación y de conciliación}



\subanexo{C.1 Ficha de observación cronometrada}

\vspace{8pt}{\footnotesize\setstretch{1.0}
\begin{longtable}{>{\raggedright\arraybackslash}p{0.1150\dimexpr\linewidth-14\tabcolsep\relax}>{\raggedright\arraybackslash}p{0.2414\dimexpr\linewidth-14\tabcolsep\relax}>{\raggedright\arraybackslash}p{0.0920\dimexpr\linewidth-14\tabcolsep\relax}>{\raggedright\arraybackslash}p{0.0920\dimexpr\linewidth-14\tabcolsep\relax}>{\raggedright\arraybackslash}p{0.1034\dimexpr\linewidth-14\tabcolsep\relax}>{\raggedright\arraybackslash}p{0.1724\dimexpr\linewidth-14\tabcolsep\relax}>{\raggedright\arraybackslash}p{0.1839\dimexpr\linewidth-14\tabcolsep\relax}}
\caption{Registro de tiempos de los procesos observados}\label{tab:v4_C1}\\
\toprule
\textbf{Fecha} & \textbf{Actividad} & \textbf{Inicio} & \textbf{Fin} & \textbf{Duración} & \textbf{Errores o reprocesos} & \textbf{Observaciones} \\
\midrule\endfirsthead
\toprule
\textbf{Fecha} & \textbf{Actividad} & \textbf{Inicio} & \textbf{Fin} & \textbf{Duración} & \textbf{Errores o reprocesos} & \textbf{Observaciones} \\
\midrule\endhead
\bottomrule\endfoot
\bottomrule\endlastfoot
 & Registro de servicio / Cierre de caja / Liquidación &  &  &  &  &  \\
 & Registro de servicio / Cierre de caja / Liquidación &  &  &  &  &  \\
 & Registro de servicio / Cierre de caja / Liquidación &  &  &  &  &  \\
\end{longtable}}

\subanexo{C.2 Ficha de conciliación de cálculos}

\vspace{8pt}{\small\setstretch{1.0}
\begin{longtable}{>{\raggedright\arraybackslash}p{0.1905\dimexpr\linewidth-10\tabcolsep\relax}>{\raggedright\arraybackslash}p{0.2262\dimexpr\linewidth-10\tabcolsep\relax}>{\raggedright\arraybackslash}p{0.2262\dimexpr\linewidth-10\tabcolsep\relax}>{\raggedright\arraybackslash}p{0.1786\dimexpr\linewidth-10\tabcolsep\relax}>{\raggedright\arraybackslash}p{0.1786\dimexpr\linewidth-10\tabcolsep\relax}}
\caption{Comparación de comisiones calculadas por el sistema y por el cálculo de referencia}\label{tab:v4_C2}\\
\toprule
\textbf{Registro} & \textbf{Comisión del sistema (Bs)} & \textbf{Comisión de referencia (Bs)} & \textbf{Diferencia (Bs)} & \textbf{Coincide (sí/no)} \\
\midrule\endfirsthead
\toprule
\textbf{Registro} & \textbf{Comisión del sistema (Bs)} & \textbf{Comisión de referencia (Bs)} & \textbf{Diferencia (Bs)} & \textbf{Coincide (sí/no)} \\
\midrule\endhead
\bottomrule\endfoot
\bottomrule\endlastfoot
 &  &  &  &  \\
 &  &  &  &  \\
 &  &  &  &  \\
\end{longtable}}

\anexotabla{D}

\anexo{Anexo D. Guía de entrevista semiestructurada}

Entrevista dirigida a la administradora, con una duración aproximada de 30 a 40 minutos. Se grabará únicamente con su consentimiento (Anexo E).


\pregunta{\textbf{1.} ¿Cómo describe el funcionamiento diario del salón, desde que llega una clienta hasta que se cobra el servicio?}

\pregunta{\textbf{2.} ¿Qué parte del proceso le toma más tiempo o le genera más preocupación?}

\pregunta{\textbf{3.} ¿Qué situaciones le han generado errores o diferencias en la caja o en las comisiones?}

\pregunta{\textbf{4.} ¿Cómo calcula y paga la liquidación semanal? ¿Qué parte le resulta más difícil?}

\pregunta{\textbf{5.} ¿Qué opina de registrar los servicios desde el celular durante la jornada?}

\pregunta{\textbf{6.} ¿Qué expectativas tiene sobre un sistema que reemplace la planilla?}

\pregunta{\textbf{7.} ¿Qué condiciones considera indispensables para dejar de usar la planilla?}

\pregunta{\textbf{8.} ¿Qué aspecto de la planilla actual no debería perderse en el nuevo sistema?}

\anexotabla{E}

\anexo{Anexo E. Formato de consentimiento informado}

Usted está invitada a participar en el trabajo final de investigación «Sistema Web Móvil para la Gestión de Caja y Nómina del Salón de Belleza Lush Beauty». Su participación consiste en responder un cuestionario, participar en una entrevista y, si corresponde, usar el sistema durante el piloto.


La participación es voluntaria y puede retirarse en cualquier momento, sin consecuencias. La información se utilizará solo con fines académicos. Los datos de sueldos y de clientes se tratarán de forma confidencial y se presentarán de manera agregada o con códigos.


Nombre: \blank{30}     Firma: \blank{22}     Fecha: \blank{4}/\blank{4}/\blank{8}


\anexotabla{F}

\anexo{Anexo F. Planilla de ejemplo anonimizada}

Captura de la planilla de ejemplo de agosto de 2026, con los nombres de clientes sustituidos por códigos. [Insertar imagen anonimizada.]


\anexotabla{G}

\anexo{Anexo G. Ejemplo de cuestionario de requerimientos respondido}

\textbf{Ejemplo ilustrativo.} Las respuestas siguientes no provienen de una aplicación real del cuestionario. Se construyeron a partir de la información disponible en el documento de definición del producto (PDR) del proyecto y de la planilla de ejemplo de agosto de 2026. Las respuestas marcadas como «[por confirmar]» deben completarse con la administradora antes de utilizar el cuestionario como fuente de datos.


\vspace{8pt}{\small\setstretch{1.0}
\begin{longtable}{>{\raggedright\arraybackslash}p{0.0692\dimexpr\linewidth-6\tabcolsep\relax}>{\raggedright\arraybackslash}p{0.3522\dimexpr\linewidth-6\tabcolsep\relax}>{\raggedright\arraybackslash}p{0.5786\dimexpr\linewidth-6\tabcolsep\relax}}
\caption{Respuestas ilustrativas al cuestionario de requerimientos}\label{tab:v4_G1}\\
\toprule
\textbf{N.º} & \textbf{Pregunta (resumida)} & \textbf{Respuesta ilustrativa} \\
\midrule\endfirsthead
\toprule
\textbf{N.º} & \textbf{Pregunta (resumida)} & \textbf{Respuesta ilustrativa} \\
\midrule\endhead
\bottomrule\endfoot
\bottomrule\endlastfoot
1 & Rol en el salón & Ambos (administración y especialista) \\
2 & Años operando el salón & [por confirmar] \\
3 & Personas que prestan servicios & 3 \\
4 & Servicios en un día habitual & Aprox. 3 (91 servicios en agosto de 2026, según la planilla de ejemplo) \\
5 & Dispositivos para registrar & Celular; hoja de cálculo \\
6 & Forma de registro actual & Hoja de cálculo \\
7 & Momento del registro & Al final del día [por confirmar] \\
8 & Datos que registra & Fecha, cliente, especialista, servicio, precio, descuento, propina y observaciones; no registra método de pago \\
9 & Métodos de pago & Efectivo, QR, transferencia y tarjeta [por confirmar] \\
10 & Cobros combinados (combos) & Sí, a veces \\
11 & Cuadre de caja al cierre & Suma manual de montos; no hay desglose por método de pago \\
12 & Descuentos y autorización & Porcentaje fijo (0 a 50~\%); autoriza la administradora \\
13 & Servicio cancelado & [por confirmar] \\
14 & Porcentaje de comisión & 35~\% \\
15 & Base de cálculo de la comisión & Precio de lista \\
16 & La comisión varía & No \\
17 & Explicación & No aplica \\
18 & Frecuencia de pago al personal & Semanal \\
19 & Sueldo fijo además de la comisión & [por confirmar] \\
20 & Deducciones y cálculo & Adelantos, faltas y materiales; monto definido a mano [por confirmar] \\
21 & Deducciones mayores a lo ganado & [por confirmar] \\
22 & Forma de entregar el pago y el detalle & [por confirmar] \\
23 & Registro de propinas & En una columna aparte, sin pasar por la comisión \\
24 & Reparto de propinas & 100~\% a quien atendió [por confirmar] \\
25 & Necesidad de editar propinas & Sí \\
26 & Registro de clientes & Hoja de cálculo, como texto libre \\
27 & Datos de cada clienta & Nombre, teléfono y notas \\
28 & Historial de servicios por clienta & Sí \\
29 & Nombres duplicados o confusos & Sí: nombres con anotaciones y espacios finales \\
30 & Frecuencia de cambio de precios & [por confirmar] \\
31 & Precios que varían al cobro & Sí: diseños (base más Bs~20) y arreglos sin precio publicado \\
32 & Quién modifica catálogo y precios & La administradora \\
33 & Reportes necesarios & Caja diaria, detalle por servicio y liquidación por persona \\
34 & Formato de los reportes & Excel, y PDF para compartir \\
35 & Migrar la planilla actual & Sí \\
36 & Quién recibe reportes fuera del salón & Contador \\
37 & Frecuencia de corte de internet & [por confirmar] \\
38 & Registrar sin internet & Sí \\
39 & Quién ve sueldos y comisiones & Solo la administradora \\
40 & Quién modifica o elimina registros & Solo la administradora \\
41 & Bitácora de cambios & Sí \\
42 & Priorización (1 a 5) & Registro rápido 5; comisión automática 5; liquidación con deducciones 5; PDF para WhatsApp 4; sin conexión 4; importar Excel 4; bitácora 4; reportes para el contador 3; clientes con historial 3; agenda en línea 1 \\
43 & Problema que más preocupa & Errores de cálculo y tiempo de liquidación (ilustrativo) \\
44 & Parte del proceso que no debe cambiar & [por confirmar] \\
\end{longtable}}

\anexotabla{H}

\anexo{Anexo H. Matriz inicial de requisitos funcionales}

Los requisitos siguientes se presentan como propuesta de implementación, y deben validarse mediante entrevistas y revisión técnica antes de convertirse en línea base aprobada.


\vspace{8pt}{\small\setstretch{1.0}
\begin{longtable}{>{\raggedright\arraybackslash}p{0.3333\dimexpr\linewidth-6\tabcolsep\relax}>{\raggedright\arraybackslash}p{0.3333\dimexpr\linewidth-6\tabcolsep\relax}>{\raggedright\arraybackslash}p{0.3333\dimexpr\linewidth-6\tabcolsep\relax}}
\caption{Priorización sugerida; no sustituye las decisiones de las participantes del salón}\label{tab:dss_A1}\\
\toprule
\textbf{Código} & \textbf{Requisito propuesto} & \textbf{Prioridad} \\
\midrule\endfirsthead
\toprule
\textbf{Código} & \textbf{Requisito propuesto} & \textbf{Prioridad} \\
\midrule\endhead
\bottomrule\endfoot
\bottomrule\endlastfoot
RF-01 & Registrar venta con una o varias líneas & Crítica \\
RF-02 & Asignar especialista y servicio del catálogo & Crítica \\
RF-03 & Registrar pago por efectivo, QR o transferencia & Crítica \\
RF-04 & Aplicar descuentos con motivo y aprobación & Crítica \\
RF-05 & Calcular comisión según regla versionada & Crítica \\
RF-06 & Registrar propinas y deducciones & Crítica \\
RF-07 & Generar y aprobar liquidación semanal & Crítica \\
RF-08 & Cerrar y conciliar caja diaria & Crítica \\
RF-09 & Generar comprobantes PDF y exportaciones XLSX & Alta \\
RF-10 & Registrar servicio sin conectividad & Alta \\
RF-11 & Sincronizar idempotentemente operaciones & Crítica \\
RF-12 & Importar y mapear hojas XLSX & Alta \\
RF-13 & Proponer aliases y detectar duplicados & Alta \\
RF-14 & Registrar correcciones y fuentes de importación & Alta \\
RF-15 & Consultar dashboard financiero & Alta \\
RF-16 & Consultar cinco intenciones de Lush Insight & Alta \\
RF-17 & Detectar alertas de reglas Lush Guard & Alta \\
RF-18 & Automatizar tres eventos aprobados con n8n & Alta \\
RF-19 & Aislar permisos y datos de cada salón & Crítica \\
RF-20 & Registrar auditoría y cambios históricos & Crítica \\
\end{longtable}}

\anexo{Anexo I. Historias de usuario para el backlog inicial}

\textbf{HU-01. Como administradora, quiero registrar una venta con servicios y métodos de pago, para gestionar con trazabilidad el proceso.}


\textit{Criterio de aceptación:} el importe de líneas coincide con pagos y se genera detalle.


\textbf{HU-02. Como especialista, quiero registrar un servicio desde el móvil sin internet, para gestionar con trazabilidad el proceso.}


\textit{Criterio de aceptación:} la cola muestra pendiente y luego confirmada sin duplicación.


\textbf{HU-03. Como administradora, quiero establecer una política de comisiones vigente, para gestionar con trazabilidad el proceso.}


\textit{Criterio de aceptación:} el cálculo usa un snapshot histórico y conserva motivo de cambios.


\textbf{HU-04. Como administradora, quiero liquidar una semana con propinas y deducciones, para gestionar con trazabilidad el proceso.}


\textit{Criterio de aceptación:} se obtiene un total verificable y un PDF por especialista.


\textbf{HU-05. Como administradora, quiero importar mi planilla anterior, para gestionar con trazabilidad el proceso.}


\textit{Criterio de aceptación:} recibo un reporte de accepted/reviewed/rejected antes de confirmar.


\textbf{HU-06. Como administradora, quiero consultar por qué cambió el neto, para gestionar con trazabilidad el proceso.}


\textit{Criterio de aceptación:} el asistente consulta el período y muestra la fuente del cálculo.


\textbf{HU-07. Como administradora, quiero recibir una alerta si se supera un umbral, para gestionar con trazabilidad el proceso.}


\textit{Criterio de aceptación:} la alerta indica regla y permite marcar revisión.


\textbf{HU-08. Como administradora, quiero enviar un aviso al aprobar caja, para gestionar con trazabilidad el proceso.}


\textit{Criterio de aceptación:} el flujo n8n registra éxito o fallo y permite reintentar.


\textbf{HU-09. Como especialista, quiero ver mi propia liquidación, para gestionar con trazabilidad el proceso.}


\textit{Criterio de aceptación:} no puedo consultar las de otras especialistas.


\textbf{HU-10. Como administradora, quiero revertir una venta con motivo, para gestionar con trazabilidad el proceso.}


\textit{Criterio de aceptación:} se conserva historial y se genera ajuste correspondiente.


\anexotabla{J}

\anexo{Anexo J. Casos de prueba de aceptación}

\vspace{8pt}{\small\setstretch{1.0}
\begin{longtable}{>{\raggedright\arraybackslash}p{0.3333\dimexpr\linewidth-6\tabcolsep\relax}>{\raggedright\arraybackslash}p{0.3333\dimexpr\linewidth-6\tabcolsep\relax}>{\raggedright\arraybackslash}p{0.3333\dimexpr\linewidth-6\tabcolsep\relax}}
\caption{Pruebas propuestas; usar datos ficticios y ejecutar todas antes de la feria}\label{tab:dss_C1}\\
\toprule
\textbf{Caso} & \textbf{Preparación / acción} & \textbf{Resultado esperado} \\
\midrule\endfirsthead
\toprule
\textbf{Caso} & \textbf{Preparación / acción} & \textbf{Resultado esperado} \\
\midrule\endhead
\bottomrule\endfoot
\bottomrule\endlastfoot
CP-01 & Venta Bs~120, descuento autorizado Bs~8 & Se conserva lista, descuento y cobro neto Bs~112 \\
CP-02 & Comisión 35~\% con base lista sobre Bs~120 & Comisión de Bs~42,00 con regla histórica \\
CP-03 & Reintentar dos veces UUID idéntico & Solo una venta confirmada \\
CP-04 & Cortar internet antes de registrar & Registro local pendiente y posterior confirmación \\
CP-05 & Importar “ PEDI” y “PEDI” & Sugerencia de un mismo servicio antes de aprobar \\
CP-06 & Especialista consulta nómina ajena & Acceso rechazado y registro de incidencia \\
CP-07 & IA responde una cifra sin datos del período & Indica falta de información, no inventa valor \\
CP-08 & Emitir dos veces evento cierre\_caja.aprobado & Un solo aviso lógico para el mismo evento \\
CP-09 & Modificar una política ya aplicada & No se recalculan ventas históricas \\
CP-10 & Error del correo al enviar un PDF & Nómina persiste; notificación se reintenta \\
\end{longtable}}

\anexo{Anexo K. Diseño inicial de API y eventos}

\textbf{Endpoints propuestos}


\begin{vinetas}
\item \texttt{POST /api/v1/ventas/} — crea una venta idempotente.
\item \texttt{GET /api/v1/caja/resumen/?desde=\&hasta=} — consolida cobros y descuentos.
\item \texttt{POST /api/v1/caja/cierres/} — registra cierre y diferencias.
\item \texttt{POST /api/v1/liquidaciones/generar/} — crea borrador para un período.
\item \texttt{POST /api/v1/liquidaciones/\{id\}/aprobar/} — autoriza una liquidación.
\item \texttt{POST /api/v1/importaciones/} — carga y mapea un XLSX.
\item \texttt{POST /api/v1/importaciones/\{id\}/confirmar/} — aplica correcciones aprobadas.
\item \texttt{GET /api/v1/alertas/} — consulta incidencias según permisos.
\item \texttt{POST /api/v1/insight/preguntar/} — interpreta una pregunta de lectura.
\item \texttt{GET /api/v1/auditoria/} — consulta bitácora según rol.
\end{vinetas}

\textbf{Eventos de integración}


\begin{vinetas}
\item caja.cierre\_aprobado, con event\_id, salon\_id, cierre\_id, timestamp y versión.
\item nomina.liquidacion\_aprobada, con event\_id, salon\_id, liquidacion\_id y versión.
\item control.alerta\_pendiente, con event\_id, salon\_id, alerta\_id y severidad de negocio.
\end{vinetas}

Los contratos finales incluirán autenticación, códigos HTTP, validación de datos, paginación, límites de tasa y pruebas de compatibilidad. Los eventos no transmitirán nombres completos ni importes personales si no son necesarios.


\anexo{Anexo L. Diseño operativo de los flujos n8n}

\textbf{E1. Cierre de caja}


\begin{vinetas}
\item Disparador: evento firmado caja.cierre\_aprobado.
\item Validación: event\_id único y esquema permitido.
\item Consulta: GET de resumen con cuenta de servicio de solo lectura.
\item Salida: correo al responsable con período y totales.
\item Contingencia: persistir error y reintentar sin modificar caja.
\end{vinetas}

\textbf{E2. Liquidación}


\begin{vinetas}
\item Disparador: evento nomina.liquidacion\_aprobada.
\item Validación: autorización de destinatario y vigencia del enlace.
\item Salida: aviso con acceso restringido al PDF individual.
\item Contingencia: correo fallido no invalida la liquidación aprobada.
\end{vinetas}

\textbf{E3. Alerta}


\begin{vinetas}
\item Disparador: control.alerta\_pendiente.
\item Validación: incidencia no notificada anteriormente.
\item Salida: mensaje con resumen no sensible y enlace de revisión.
\item Contingencia: registrar error y escalar al panel de incidencias.
\end{vinetas}

\anexo{Anexo M. Instrumentos adicionales para el piloto}

La extensión del cuestionario de requerimientos incorporará preguntas específicas para decidir qué funciones son útiles, evitando asumir necesidades que no hayan sido confirmadas por las usuarias.


\begin{vinetas}
\item ¿Qué información se necesita consultar al finalizar cada jornada?
\item ¿Qué diferencias de caja o errores de nombres son más frecuentes?
\item ¿De qué forma se revisan y autorizan descuentos especiales?
\item ¿Con qué frecuencia se pierde conectividad durante el registro?
\item ¿Qué datos de nómina puede visualizar cada especialista?
\item ¿Qué tipo de aviso automático sería útil, y por qué canal autorizado?
\item ¿Qué preguntas sobre resultados económicos sería importante responder?
\item ¿Qué sugerencias automáticas necesitan confirmación antes de aplicarse?
\end{vinetas}

\anexo{Anexo N. Lista de preparación para la feria}

\begin{vinetas}
\item[$\square$] Datos ficticios anonimizados y separados del entorno de piloto.
\item[$\square$] Teléfono con PWA instalada y cola local previamente verificada.
\item[$\square$] Archivo de prueba XLSX que contiene nombres inconsistentes conocidos.
\item[$\square$] Reglas financieras revisadas y aprobadas por la administradora.
\item[$\square$] Cuenta de demostración sin permisos a datos reales de nómina.
\item[$\square$] Dos preguntas de Lush Insight ensayadas y reporte fuente disponible.
\item[$\square$] Flujos n8n visibles con registro de ejecución y reintentos controlados.
\item[$\square$] Respaldos, entorno local y video de contingencia preparados.
\item[$\square$] Guion de cinco minutos con responsables de cada segmento.
\item[$\square$] Métricas de pruebas impresas sin promesas de resultados reales.
\end{vinetas}

\anexo{Anexo O. Fichas de decisión para el caso de feria}

\textbf{Ficha DEC-01.} Problema: decidir entre mantener precio, descuento del 5~\% y descuento del 10~\% para un servicio. Propietaria: administradora. Datos observados: precio de lista y política vigente, después de validación. Datos supuestos: costos variables y volumen de ventas. Métrica primaria: contribución agregada. Resultado esperado: una alternativa comparada y justificada.


\textbf{Ficha DEC-02.} Sensibilidad: modificar el volumen estimado entre 90 y 140 servicios y variar el peso del criterio económico. Resultado esperado: se identifica si cambia la alternativa preferida, y se explica por qué.


\textbf{Ficha DEC-03.} Seguimiento: fecha de evaluación posterior, indicador esperado y observado, observaciones de factores externos y decisión de mantener o revisar la política.


\anexo{Anexo P. Casos de prueba de soporte para decisiones}

\textbf{DSS-CP01.} Precio Bs~120, descuento 10~\%, comisión 35~\% sobre lista y materiales Bs~20: precio cobrado Bs~108, comisión Bs~42, contribución unitaria Bs~46.


\textbf{DSS-CP02.} Para 100 servicios sin descuento: contribución total Bs~5.800. Para descuento 10~\%, volumen mínimo entero para igualar el total: 127 servicios.


\textbf{DSS-CP03.} Pesos de matriz suman 100~\%. Un total distinto debe bloquear la evaluación o normalizarse únicamente con confirmación visible.


\textbf{DSS-CP04.} Un escenario guardado no modifica registros de ventas, remuneraciones ni políticas de comisión.


\textbf{DSS-CP05.} Una recomendación con costos faltantes indica insuficiencia de datos y no informa utilidad neta.


\textbf{DSS-CP06.} Un usuario sin rol de administración no aprueba cambios de precios ni accede a nómina ajena.


\textbf{DSS-CP07.} El historial conserva resultado esperado, responsable y revisión posterior; no reescribe el escenario original.


\anexo{Anexo Q. Instrumento de evaluación de decisiones}

\textbf{Ficha por tarea:} ID de participante anonimizado; fecha; modalidad (Excel/manual o DSS); descripción de la tarea; datos entregados; hora de inicio y fin; solución producida; criterio usado; fuentes citadas; advertencia de supuestos; autorización humana; comentarios de usabilidad.


\textbf{Rúbrica propuesta (0 a 2 puntos por criterio):} (1) cálculo correcto, (2) alternativas comparadas, (3) supuestos explícitos, (4) explicación comprensible y (5) decisión trazable. Un puntaje alto indica desempeño en la tarea, no demuestra automáticamente mejores resultados empresariales.


\textbf{Cuestionario de interpretación:} ¿qué datos provienen de operaciones reales?, ¿qué valores son supuestos?, ¿qué criterio influyó más en la recomendación?, ¿qué cambio podría invertir la decisión?, ¿quién debe aprobar la acción?

